\section{Криптографические основы многосторонних конфиденциальных вычислений}
\label{sec:cryptographic-foundations}

\subsection{Разделение секрета}

\subsubsection{Определение и свойства}

Разделение секрета --- это криптографическая схема, позволяющая разбить секрет \(S\) на \(n\) частей (долей) \(S_1,\ldots,S_n\) так, что любой набор не менее \(k\) долей позволяет восстановить \(S\), а из меньшего числа долей невозможно получить какую-либо информацию о секрете. Для такой \((k,n)\)-пороговой схемы характерны два свойства: во-первых, знание любых \(k\) или более долей позволяет восстановить секрет; во-вторых, знание не более \(k-1\) долей не уменьшает неопределённость секрета. Схема считается идеальной и совершенной, если доли имеют ту же длину, что и секрет, и обладают информационно-теоретической стойкостью.

\subsubsection{Схема разделения секрета Шамира}

Наиболее известной реализацией является схема Шамира (Shamir’s secret sharing).  Пусть секрет \(S\) отображается на элемент \(a_0\) конечного поля \(\mathrm{GF}(q)\).  Для порогового значения \(k\) случайно выбирается полином степени \(k-1\)

\begin{equation}
f(x) = a_0 + a_1x + a_2x^2 + \cdots + a_{k-1}x^{k-1},
\end{equation}

где \(a_1,\ldots,a_{k-1}\) --- случайные коэффициенты из \(\mathrm{GF}(q)\).  Для \(i=1,\ldots,n\) вычисляются значения \(f(i)\) и образуют пары \((i, f(i))\); каждому участнику выдаётся одна такая пара.  Любой набор из \(k\) пар позволяет восстановить полином \(f\) интерполяцией Лагранжа:

\begin{equation}
f(x) = \sum_{j=0}^{k-1} y_j \ell_j(x), \qquad
\ell_j(x) = \prod_{\substack{m=0\\m \neq j}}^{k-1} \frac{x-x_m}{x_j-x_m}.
\label{eq:lagrange-interpolation}
\end{equation}

Здесь \((x_j,y_j)\) --- выбранные доли, где \(y_j=f(x_j)\), а \(\ell_j(x)\) --- базисные многочлены, равные единице при \(x=x_j\) и нулю в остальных выбранных точках. Поэтому формула~\eqref{eq:lagrange-interpolation} однозначно восстанавливает полином степени не выше \(k-1\) по \(k\) различным точкам. Подстановка \(x=0\), для которой \(f(0)=a_0\), даёт формулу восстановления секрета:

\begin{equation}
a_0 = f(0) = \sum_{j=0}^{k-1} y_j \prod_{\substack{m=0\\m \neq j}}^{k-1} \frac{x_m}{x_m - x_j},
\end{equation}

\subsubsection{Пример}

Для наглядности рассмотрим вычисления в поле \(\mathrm{GF}(65521)\), где 65521 --- простое число, а сложение, умножение и деление выполняются по модулю 65521. Секрету соответствует элемент поля \(a_0=1234\). В схеме \((3,6)\) выбирается случайный полином степени не выше двух. Пусть в одном из запусков случайные коэффициенты приняли значения \(a_1=166\) и \(a_2=94\). Тогда полином имеет вид

\begin{equation}
f(x) = 1234 + 166x + 94x^2 \pmod{65521}.
\end{equation}

Каждому из шести участников передаётся пара \((i,f(i))\), где \(i\) --- его уникальный ненулевой номер в поле. Полученные доли равны:

\begin{equation}
\begin{aligned}
f(1) &= 1494, & f(2) &= 1942, & f(3) &= 2578,\\
f(4) &= 3402, & f(5) &= 4414, & f(6) &= 5614.
\end{aligned}
\end{equation}

Все значения в примере меньше 65521, поэтому их можно записать без приведения по модулю; в общем случае участники хранят элементы поля. Для восстановления достаточно собрать любые три доли с различными номерами. Например, для долей \((1,1494)\), \((2,1942)\) и \((3,2578)\) коэффициенты интерполяции Лагранжа в точке \(x=0\) согласно формуле~\eqref{eq:lagrange-interpolation} равны:

\begin{equation}
\begin{aligned}
\lambda_1 &= \frac{(0-2)(0-3)}{(1-2)(1-3)}=3, \\
\lambda_2 &= \frac{(0-1)(0-3)}{(2-1)(2-3)}=-3, \\
\lambda_3 &= \frac{(0-1)(0-2)}{(3-1)(3-2)}=1.
\end{aligned}
\end{equation}

Следовательно, секрет восстанавливается как свободный член полинома:

\begin{equation}
f(0) \equiv 3f(1)-3f(2)+f(3)
\equiv 3\cdot1494-3\cdot1942+2578
\equiv 1234 \pmod{65521}.
\end{equation}

Та же процедура с любой другой тройкой долей даст тот же результат, поскольку три точки с различными номерами однозначно задают полином степени не выше двух. Одна или две доли, напротив, не позволяют отличить один секрет от другого: для каждого возможного значения \(a_0\) можно подобрать случайные коэффициенты, согласующиеся с полученными долями. При равномерном независимом выборе коэффициентов распределение наблюдаемых долей поэтому не зависит от секрета.

В практическом применении поле выбирают так, чтобы в нём хватало различных ненулевых номеров участников, а секрет кодировался элементом этого поля; более длинные секреты при необходимости представляют несколькими элементами. Для каждого нового распределения долей коэффициенты выбирают заново. Базовая схема обеспечивает конфиденциальность долей, но не проверяет их подлинность: ошибочная или намеренно изменённая доля может привести к неверному результату, поэтому при активном противнике нужны дополнительные механизмы проверки.

\subsubsection{Аддитивное и гомоморфное деление секрета}

Помимо схемы Шамира используется простое аддитивное разделение, при котором секрет разбивается на \(n\) случайных элементов поля так, что их сумма равна \(S\).  Например, для секрета \(s\) выбираются \(n-1\) случайных чисел \(r_1,\ldots,r_{n-1}\) и вычисляется \(r_n = s - \sum_{i=1}^{n-1} r_i\).  В дальнейшем операции над секретами превращаются в поэлементные операции над долями: суммируя доли по модулю \(q\), участники получают долю суммы исходных секретов.  Этот подход лежит в основе протоколов SPDZ и BMR, где для мультипликативных операций используются так называемые тройки Бивера (Beaver triples) \cite{06}.  Гомоморфное деление секрета (homomorphic secret sharing) сочетает идеи деления и шифрования и позволяет выполнять локальные операции над долями без дальнейшей коммуникации.

\subsection{Искаженные схемы (garbled circuits)}

Зашифрованная схема --- криптографический протокол, позволяющий двум сторонам вычислять любую булеву функцию от своих входных данных без раскрытия этих данных.  Протокол был предложен Яо и требует, чтобы вычисляемая функция была представлена в виде булевой схемы с двумя входами.

Пусть Алиса (garbler) и Боб (evaluator) хотят вычислить значение функции \(f(x_A, x_B)\).  Протокол состоит из следующих шагов:

\begin{itemize}
\item Построение схемы: участники согласуют общую булеву схему (набор логических элементов И, ИЛИ и пр.), реализующую функцию \(f\).

\item Гарболизация: Алиса для каждого логического провода \(w\) схемы генерирует две случайные \(k\)-битовые метки \(X^w_0, X^w_1\) для логических значений 0 и 1. Затем она проходит по всем элементам и заменяет значения в таблице истинности соответствующими метками. После этого она шифрует каждую строку таблицы двойным симметричным шифрованием с ключами, соответствующими двум входным меткам, и случайно переставляет строки; результатом является искаженная таблица.

\item Передача схемы и меток: Алиса отправляет Бобу зашифрованные таблицы всех элементов и метки, соответствующие своим входам. Метки Боба для каждого бита он получает через протокол слепой передачи (см. п. 2.4), что гарантирует, что Алиса не узнает его выбор.

\item Оценивание: Боб, имея искаженную схему и свои метки, проходит по элементам схемы, подбирая для каждого элемента единственную строку таблицы, которую он может расшифровать. В результате он получает метку на выходном логическом проводе и декодирует её с помощью соответствующего алгоритма, получая значение функции.
\end{itemize}

В формальном описании схема искажения задается рандомизированным алгоритмом \(\text{Gb}\), который для функции \(f\) формирует тройку \((F,e,d)\): \(F\) --- искаженное представление функции, \(e\) --- функция кодирования входов, а \(d\) --- функция декодирования выхода. При вычислении \(F(e(x))\) получается искаженный выход, который после применения \(d\) преобразуется в значение \(f(x)\).

\subsection{Гомоморфное шифрование (homomorphic encryption)}

\subsubsection{Определение и свойства}

Гомоморфное шифрование --- это криптографический метод, позволяющий выполнять вычислительные операции непосредственно над зашифрованными данными без их предварительного расшифрования. При этом расшифрование результата вычислений позволяет получить значение, эквивалентное результату выполнения соответствующих операций над исходными данными. Возможность выполнения вычислений заданной глубины определяется особенностями криптографической схемы и выбранными параметрами. Для обеспечения вычислений большей глубины в ряде схем применяется процедура бутстрэппинга, позволяющая уменьшать накопленный в процессе вычислений шум в шифротекстах.

\subsubsection{Частично гомоморфное шифрование (PHE)}

В простейших схемах выполняется только одно арифметическое действие.  Например, в криптосистеме RSA (без паддинга) шифротекст \(E(m) = m^e \bmod n\) обладает мультипликативным свойством: произведение двух шифротекстов равно шифротексту произведения исходных сообщений.  Схема Эль-Гамаля использует пары \((g^r, m \cdot h^r)\), и перемножение двух шифротекстов приводит к шифрованию произведения сообщений.  В схеме Паилье (Paillier) сообщение шифруется как \(E(m) = g^m r^n \bmod n^2\), и произведение шифротекстов даёт шифр суммы: \(E(m_1)\cdot E(m_2) = E(m_1 + m_2)\).  Эти свойства позволяют строить МКВ-протоколы для простых операций без интерактивности.

\subsubsection{Полностью гомоморфное шифрование (FHE)}

Полностью гомоморфные схемы выполняют как сложение, так и умножение над шифротекстами.  Идея была предложена в 1978 году, но первый практически реализуемый FHE-протокол разработал Гентри в 2009 году, введя метод «bootstrap» для подавления шума.  Схема строится на «частично» гомоморфном шифре, который поддерживает ограниченное число операций, и периодическом гомоморфном декодировании, уменьшающем шум.  Теоретические описания FHE используют абстракцию колец и операции AND/XOR как аналоги умножения/сложения.  Существующие реализации (CKKS, BGV, BFV) основаны на сложности задач «обучения с ошибками» и «приближенного НОД»; они позволяют вычислять целые и вещественные функции над зашифрованными данными, но требуют значительных вычислительных ресурсов и больших размеров ключей.

\subsection{Протокол слепой передачи (oblivious transfer)}

Слепая передача (oblivious transfer, OT) --- семейство протоколов выборочной передачи сообщений. В распространенном варианте 1-из-2 отправитель Алиса располагает сообщениями \(m_0\) и \(m_1\), а получатель Боб задаёт приватный выбор \(b\in\{0,1\}\). По завершении протокола Боб получает сообщение \(m_b\), соответствующее его выбору. Протокол обеспечивает два свойства конфиденциальности: Алиса не узнает значение \(b\), а Боб не получает сведений о невыбранном сообщении \(m_{1-b}\), кроме тех, которые можно вывести из самого \(m_b\). Иными словами, отправитель не видит выбор получателя, а получатель узнает только выбранное сообщение.

В варианте 1-из-\(n\) отправитель имеет сообщения \(m_1,\ldots,m_n\), а получатель выбирает индекс \(i\) и получает только \(m_i\), не раскрывая отправителю свой выбор. Исторически ранний вариант Рабина устроен иначе: отправитель передает одно сообщение, которое получатель узнает с вероятностью \(1/2\), при этом отправитель не знает, состоялась ли передача. Этот вариант не следует смешивать с 1-из-2 OT, где получатель сам выбирает одно из двух сообщений \cite{06}.

OT является полным примитивом для безопасных вычислений: с его помощью и локальными вычислениями можно построить протокол для любой полиномиально вычислимой функции в рамках выбранной модели безопасности \cite{06}.


Приведем пример протокола 1-из-2 на основе RSA: Алиса (отправитель) хочет передать Бобу (получателю) одно из сообщений \(m_0\) и \(m_1\), не узнавая, какое именно сообщение выбрал Боб, и не раскрывая ему второе сообщение. Для простоты каждое сообщение рассматривается как элемент кольца вычетов \(\mathbb{Z}_N\).

Сначала Алиса генерирует ключ RSA: выбирает простые числа \(p\) и \(q\), вычисляет \(N=pq\) и \(\varphi(N)=(p-1)(q-1)\), затем выбирает \(e\), взаимно простое с \(\varphi(N)\), и находит \(d\), для которого \(ed\equiv1\pmod{\varphi(N)}\). Пара \((N,e)\) является открытым ключом, а \(d\) Алиса сохраняет в тайне. Она также выбирает два различных случайных значения \(x_0,x_1\in\mathbb{Z}_N\) и сообщает их Бобу.

Боб выбирает бит \(b\in\{0,1\}\), обозначающий нужное сообщение, и случайное \(k\in\mathbb{Z}_N^*\). Он вычисляет значение \(v\) и отправляет его Алисе. Получив \(v\), Алиса вычисляет для каждого \(i\in\{0,1\}\) маску \(r_i\), складывает её с соответствующим сообщением по модулю \(N\) и отправляет Бобу оба результата:

\begin{equation}
\begin{aligned}
v &\equiv x_b + k^e \pmod N, \\
r_i &\equiv (v-x_i)^d \pmod N, \quad i\in\{0,1\}, \\
m'_i &\equiv m_i + r_i \pmod N, \quad i\in\{0,1\}, \\
m_b &\equiv m'_b-k \pmod N.
\end{aligned}
\end{equation}

Здесь \(x_0,x_1\) --- случайные значения Алисы, связывающие вычисления с двумя вариантами; \(b\) --- выбор Боба; \(k\) --- его секретное случайное значение; \(v\) --- передаваемый Алисе запрос; \(r_i\) --- маска, полученная с помощью закрытого показателя RSA; \(m'_i\) --- замаскированное сообщение. Так как \(v-x_b\equiv k^e\pmod N\), для выбранного индекса \(r_b\equiv(k^e)^d\equiv k\pmod N\). Поэтому Боб вычитает известное ему \(k\) из \(m'_b\) и получает \(m_b\). Для второго индекса он не знает соответствующую маску \(r_{1-b}\): ее вычисление требует обращения RSA-преобразования. Алиса вычисляет обе маски, но при равномерном выборе \(x_0,x_1\) и достаточно большом модуле значение \(v\) практически не раскрывает ей выбор \(b\). Эти свойства относятся к базовой схеме при корректном выполнении протокола и опираются на вычислительную трудность обращения RSA.

Проверим вычисления на маленьком примере, который нужен только для иллюстрации и не обеспечивает криптографическую защиту. Пусть \(p=5\), \(q=11\), тогда \(N=55\) и \(\varphi(N)=40\). Выберем \(e=3\), \(d=27\), поскольку \(3\cdot27\equiv1\pmod{40}\). Пусть Алиса задала \(x_0=7\), \(x_1=12\) и сообщения \(m_0=20\), \(m_1=31\). Боб выбирает \(b=1\) и \(k=4\). Тогда:

\begin{equation}
\begin{aligned}
v &\equiv 12+4^3 \equiv 21 \pmod{55}, \\
r_0 &\equiv (21-7)^{27} \equiv 14^{27} \equiv 9 \pmod{55}, \\
r_1 &\equiv (21-12)^{27} \equiv 9^{27} \equiv 4 \pmod{55}, \\
m'_0 &\equiv 20+9 \equiv 29 \pmod{55}, \\
m'_1 &\equiv 31+4 \equiv 35 \pmod{55}.
\end{aligned}
\end{equation}

Алиса отправляет Бобу пару \((29,35)\). Поскольку Боб выбрал индекс \(b=1\), он использует второе значение и вычисляет \(m_1\equiv35-4\equiv31\pmod{55}\). Поскольку модуль \(N=55\) легко факторизовать, Боб в этом примере может вычислить и маску \(r_0\); расчет иллюстрирует только восстановление выбранного сообщения, но не обеспечивает конфиденциальность второго. В реальном применении используют большие ключи и специально анализируемые протоколы; приведенные малые параметры и учебная схема не подходят для защиты данных от активного злоумышленника.


\subsection{Сравнительный анализ криптографических примитивов}

В таблице~\ref{tab:mpc-primitives} сопоставлены криптографические примитивы МКВ; выбор зависит от модели угроз, числа участников, вычисляемой функции и допустимых затрат.

{\small
\setlength{\tabcolsep}{3pt}
\renewcommand{\arraystretch}{1.2}
\setlength{\tablecontentwidth}{\dimexpr\textwidth-10\tabcolsep-6\arrayrulewidth\relax}
\begin{longtable}{|P{0.16\tablecontentwidth}|P{0.20\tablecontentwidth}|P{0.22\tablecontentwidth}|P{0.22\tablecontentwidth}|P{0.20\tablecontentwidth}|}
\caption{Сравнение криптографических примитивов МКВ}\label{tab:mpc-primitives}\\
\hline
\multicolumn{1}{|H{0.16\tablecontentwidth}|}{\textbf{Примитив}} & \multicolumn{1}{H{0.20\tablecontentwidth}|}{\textbf{Основная идея}} & \multicolumn{1}{H{0.22\tablecontentwidth}|}{\textbf{Преимущество}} & \multicolumn{1}{H{0.22\tablecontentwidth}|}{\textbf{Недостаток}} & \multicolumn{1}{H{0.20\tablecontentwidth}|}{\textbf{Область применения}} \\
\hline
\endfirsthead
\multicolumn{5}{@{}l@{}}{\normalsize Продолжение таблицы~\thetable}\\
\hline
\multicolumn{1}{|H{0.16\tablecontentwidth}|}{\textbf{Примитив}} & \multicolumn{1}{H{0.20\tablecontentwidth}|}{\textbf{Основная идея}} & \multicolumn{1}{H{0.22\tablecontentwidth}|}{\textbf{Преимущество}} & \multicolumn{1}{H{0.22\tablecontentwidth}|}{\textbf{Недостаток}} & \multicolumn{1}{H{0.20\tablecontentwidth}|}{\textbf{Область применения}} \\
\hline
\endhead
Разделение секрета & Секрет кодируется долями; \(k\) из \(n\) долей позволяют восстановить секрет. & Информационно-теоретическая стойкость; лёгкие арифметические операции; эффективная реализация для линейных функций. & Требует много раундов обмена сообщениями для умножения; каждый участник должен быть онлайн. & Протоколы SPDZ/BMR, защищённое машинное обучение, приватная аналитика данных. \\
\hline
Искаженные схемы (GC) & Функция представляется булевой схемой, её таблицы истинности шифруются случайными метками. & Небольшое число раундов; поддержка произвольных функций; удобство для сравнений. & Большие таблицы и объём передачи; обычно применяется для двух участников; криптографические операции на каждом элементе. & Задача миллионеров Яо, приватное сравнение, доказательства с нулевым разглашением. \\
\hline
Гомо\-морфное шифрование & Операции выполняются над шифротекстами; после расшифрования результат совпадает с результатом операции над открытыми данными. & Отсутствие интерактивности; одна сторона может выполнить вычисления, не раскрывая входы. & Высокие вычислительные затраты; ограниченная глубина при FHE \cite{02}; большие ключи. & Облачные вычисления, приватные ML-сервисы, вычисления на единичных данных. \\
\hline
Слепая передача (OT) & Получатель выбирает и получает одно из \(n\) сообщений, не раскрывая выбора; отправитель не узнаёт, какое сообщение выбрано. & Универсальность и полнота для МКВ; удобство для ввода скрытых данных. & Сам по себе не вычисляет функцию; требует криптографических операций, например RSA или протокола Диффи--Хеллмана. & Ввод входов в garbled circuits, PSI (Private Set Intersection), oblivious PRF. \\
\hline
\end{longtable}
}

\subsection{Вывод}

В данном разделе были рассмотрены криптографические примитивы, применяемые в МКВ, в частности разделение секрета, искаженные схемы, гомоморфное шифрование и слепая передача. Сравнение методов, в том числе их преимущества, недостатки и области применения, приведено в таблице~\ref{tab:mpc-primitives}. Выбор протокола зависит от модели угроз, задачи и допустимых затрат; эти факторы будут учтены в дальнейшем исследовании.